\documentclass[12pt]{article}

\usepackage[margin=1in]{geometry}
\usepackage{amsmath,amssymb,amsthm}
\usepackage{booktabs}
\usepackage{natbib}
\usepackage{hyperref}
\usepackage{setspace}
\usepackage[T1]{fontenc}
\usepackage{mathpazo}

\hypersetup{colorlinks=true, linkcolor=blue, citecolor=blue, urlcolor=blue}

\newtheorem{assumption}{Assumption}
\newtheorem{proposition}{Proposition}

\onehalfspacing

\title{Inferring Treatment Compliance from Delivery-Window Data}
\author{Gaurav Sood}
\date{\today}

\begin{document}

\maketitle

\begin{abstract}
\noindent In randomized experiments with imperfect compliance, the local average treatment effect (LATE) requires observing treatment receipt. I consider settings where receipt is unobserved but the experiment has a pre-treatment time series and a distinct treatment delivery window. I propose classifying compliance by testing whether each treated unit's outcome during the delivery window exceeds its pre-treatment baseline. The false positive rate of this test equals the significance level $\alpha$ by construction, producing a closed-form bias in the inferred compliance rate that can be corrected analytically. I derive the delta method variance, exploiting the independence between the compliance classifier and the outcome that arises from the temporal structure of the experiment. In simulations calibrated to a field experiment on software downloads, the corrected LATE is unbiased and the 95\% confidence interval achieves coverage between 0.947 and 0.974 across twelve parameter configurations.

\medskip
\noindent \textit{Keywords:} LATE, compliance, instrumental variables, structural break, test-and-select

\medskip
\noindent \textit{JEL codes:} C21, C26, C12
\end{abstract}

\section{Setting}

Consider a randomized experiment with $N$ units indexed by $i$. Let $Z_i \in \{0,1\}$ denote treatment assignment and $D_i \in \{0,1\}$ denote actual treatment receipt, which is unobserved. The experiment proceeds in three temporal phases: a pre-treatment period with $T_{\text{pre}}$ outcome observations per unit, a treatment delivery window with $T_{\text{treat}}$ observations during which the experimenter attempts to apply the treatment, and a post-treatment outcome period with $T_{\text{post}}$ observations from which the outcome of interest $Y_i$ is constructed.

Noncompliance is one-sided: $D_i = 0$ for all units with $Z_i = 0$. Among treated units, $D_i = 1$ with probability $\pi$ (compliers) and $D_i = 0$ with probability $1 - \pi$ (never-takers). Under the standard assumptions---independence, exclusion restriction, monotonicity, and a nonzero first stage \citep{imbens1994, angrist1996}---the LATE is identified by the Wald estimator:
\begin{equation}
\tau = \frac{E[Y_i \mid Z_i = 1] - E[Y_i \mid Z_i = 0]}{E[D_i \mid Z_i = 1] - E[D_i \mid Z_i = 0]} = \frac{\text{ITT}}{\pi}
\end{equation}
The problem is that $\pi$ is unknown because $D_i$ is unobserved.

\section{Compliance Classification}

The treatment, when successfully delivered, produces a large and deterministic shift in the unit's outcome during the delivery window. This means that compliers and never-takers have different distributions of delivery-window outcomes conditional on the pre-treatment baseline, and a statistical test can distinguish them.

For each treated unit $i$, I test whether the mean outcome during the delivery window significantly exceeds the pre-treatment mean:
\begin{equation}\label{eq:tstat}
t_i = \frac{\bar{y}_{i,\text{treat}} - \bar{y}_{i,\text{pre}}}{\hat{\sigma}_i \sqrt{1/T_{\text{pre}} + 1/T_{\text{treat}}}}
\end{equation}
where $\hat{\sigma}_i$ is the pre-treatment standard deviation of unit $i$'s daily outcome. I classify $\hat{D}_i = \mathbf{1}(p_i < \alpha)$ using a one-sided $t$-test with $T_{\text{pre}} - 1$ degrees of freedom. For control units, $\hat{D}_i = 0$ by construction.

The inferred compliance rate is $\hat{\pi} = N_{\text{treat}}^{-1} \sum_{i: Z_i=1} \hat{D}_i$. Its expectation decomposes as:
\begin{equation}\label{eq:pihat}
E[\hat{\pi}] = \pi \cdot \text{TPR} + (1 - \pi) \cdot \text{FPR}
\end{equation}
where TPR is the true positive rate and FPR is the false positive rate. Among never-takers, the delivery-window outcome has the same distribution as the pre-treatment outcome, so $\text{FPR} = \alpha$ by the definition of the test.

\begin{assumption}\label{as:tpr}
The mechanical dose during the delivery window is sufficiently large relative to organic variation that $\text{TPR} = 1$.
\end{assumption}

Under Assumption~\ref{as:tpr}, equation~\eqref{eq:pihat} simplifies to $E[\hat{\pi}] = \pi + (1-\pi)\alpha$.

\section{Corrected Estimator}

Inverting the expression for $E[\hat{\pi}]$ gives a corrected compliance rate:
\begin{equation}\label{eq:picorr}
\pi_c = \frac{\hat{\pi} - \alpha}{1 - \alpha}
\end{equation}
and a corrected LATE:
\begin{equation}\label{eq:latecorr}
\hat{\tau}_c = \frac{\widehat{\text{ITT}}}{\pi_c} = \frac{(1 - \alpha) \cdot \widehat{\text{ITT}}}{\hat{\pi} - \alpha}
\end{equation}

\begin{proposition}\label{prop:bias}
Under Assumption~\ref{as:tpr}, the bias ratio of the uncorrected Wald estimator $\widehat{\text{ITT}}/\hat{\pi}$ relative to the true LATE is $\pi / [\pi + (1-\pi)\alpha]$, and the corrected estimator $\hat{\tau}_c$ is consistent for $\tau$.
\end{proposition}

\noindent For $\pi = 0.75$ and $\alpha = 0.05$, the bias ratio is $0.75 / 0.7625 = 0.984$, a 1.6\% downward bias in the uncorrected estimator.

\section{Inference}

The corrected LATE is $\hat{\tau}_c = f(\widehat{\text{ITT}}, \hat{\pi})$ where $f(a, b) = (1-\alpha) \cdot a / (b - \alpha)$. The delta method gives:
\begin{equation}\label{eq:var}
\text{Var}(\hat{\tau}_c) \approx \left(\frac{1-\alpha}{\hat{\pi}-\alpha}\right)^{\!2} \text{Var}(\widehat{\text{ITT}}) + \left(\frac{(1-\alpha) \cdot \widehat{\text{ITT}}}{(\hat{\pi}-\alpha)^2}\right)^{\!2} \text{Var}(\hat{\pi})
\end{equation}
where $\text{Var}(\widehat{\text{ITT}})$ is the standard difference-in-means variance and $\text{Var}(\hat{\pi}) = \hat{\pi}(1-\hat{\pi})/N_{\text{treat}}$.

The two terms are independent because $\widehat{\text{ITT}}$ is computed from post-treatment outcome data and $\hat{\pi}$ is computed from delivery-window data, which are temporally non-overlapping. This independence arises from the same three-window structure that makes the compliance classification possible in the first place. It serves the same function as cross-fitting in the test-and-select estimator of \citet{hazard2023}, who propose dropping zero-compliance subgroups and using sample splitting across units for valid post-selection inference, and as the general DML framework of \citet{chernozhukov2018}. The temporal structure here provides a stronger form of separation: a structural partition of the data-generating process rather than a random partition of the sample.

The 95\% confidence interval is $\hat{\tau}_c \pm 1.96 \cdot \hat{\text{se}}$.

\section{Partial Identification When Assumption~\ref{as:tpr} Fails}

When TPR $< 1$, the corrected compliance rate $\pi_c$ overestimates $\pi$ because the correction attributes all of $\hat{\pi} - (1-\pi)\alpha$ to compliers, missing those who failed the test. The uncorrected estimator $\widehat{\text{ITT}}/\hat{\pi}$ understates the LATE because $\hat{\pi} > \pi$. Together these bound the true LATE:
\begin{equation}\label{eq:bounds}
\frac{\widehat{\text{ITT}}}{\hat{\pi}} \leq \tau \leq \frac{(1-\alpha) \cdot \widehat{\text{ITT}}}{\hat{\pi} - \alpha}
\end{equation}
The width of this interval shrinks as $\alpha \to 0$ and as compliance rises.

\section{Monte Carlo Evidence}

I simulate experiments with $N = 500$ (250 treated, 250 control), $T_{\text{pre}} = 30$, $T_{\text{treat}} = 6$, and $T_{\text{post}} = 30$. Baseline daily outcomes are drawn from $N(5, 3)$ across units with within-unit daily noise $\sigma = 2$. The mechanical dose is 17 per day during the delivery window. The causal effect is $\tau$ additional units per day for compliers during the post-treatment period, so the true LATE is $\tau \times T_{\text{post}}$. Each scenario uses 1{,}000 Monte Carlo replications.

\begin{table}[ht]
\centering
\caption{Corrected LATE: bias and 95\% CI coverage across parameter configurations.}
\label{tab:coverage}
\smallskip
\begin{tabular}{@{}lrrrrr@{}}
\toprule
Scenario & True LATE & Mean $\hat{\tau}_c$ & Bias & Coverage & CI Width \\
\midrule
Baseline ($\pi{=}.75$, $\tau{=}3$) & 90 & 89.4 & $-0.6$ & 0.974 & 42.6 \\
$\pi{=}.50$ & 90 & 89.5 & $-0.5$ & 0.972 & 65.9 \\
$\pi{=}.90$ & 90 & 89.9 & $-0.1$ & 0.959 & 33.7 \\
$\tau{=}1$ & 30 & 29.6 & $-0.4$ & 0.956 & 39.0 \\
$\tau{=}10$ & 300 & 300.1 & 0.1 & 1.000 & 73.1 \\
$\tau{=}0$ (null) & 0 & $-0.3$ & $-0.3$ & 0.947 & 38.4 \\
$N{=}100$ & 90 & 89.7 & $-0.3$ & 0.974 & 96.9 \\
$N{=}2000$ & 90 & 89.7 & $-0.3$ & 0.969 & 21.3 \\
$\sigma{=}5$ & 90 & 89.6 & $-0.4$ & 0.969 & 39.5 \\
Mechanical${=}5$/day & 90 & 90.0 & 0.0 & 0.967 & 42.8 \\
$T_{\text{pre}}{=}7$ & 90 & 89.7 & $-0.3$ & 0.963 & 42.6 \\
$\alpha{=}.01$ & 90 & 90.0 & 0.0 & 0.974 & 42.7 \\
\bottomrule
\end{tabular}
\end{table}

Table~\ref{tab:coverage} reports the results. Coverage lies between 0.947 and 0.974 in all configurations, right where a slightly conservative 95\% CI should land. Coverage holds under the null ($\tau = 0$, coverage 0.947), confirming that the CI correctly includes zero when there is no effect. CI width scales with $N^{-1/2}$ (compare $N = 100$, width 96.9, to $N = 2{,}000$, width 21.3) and with $\pi^{-1}$ (compare $\pi = 0.50$, width 65.9, to $\pi = 0.90$, width 33.7). Bootstrap percentile CIs at baseline give coverage of 0.934 with a tighter mean width of 38.0.

At baseline, the inferred compliance rate is $\hat{\pi} = 0.761$, the corrected rate is $\pi_c = 0.748$, and the true rate is $\pi = 0.750$. Sensitivity of the break test is 1.000 and specificity is 0.947, matching $1 - \alpha$.

Assumption~\ref{as:tpr} fails when the mechanical dose is small relative to organic variation. In a high-noise simulation ($\sigma = 5$, mechanical $= 5$/day), sensitivity drops to 0.823 and the LATE estimate is biased upward by approximately 20\%. At mechanical $= 10$/day with the same noise, the method recovers. The practical requirement is mechanical/$\sigma \gtrsim 2$.

\section{Discussion}

The method developed here exploits a feature specific to experiments with repeated pre-treatment measurements, a distinct delivery window, and a mechanical dose large enough to produce a detectable break. The approach connects to several recent threads. \citet{hazard2023} propose a test-and-select LATE estimator that drops zero-compliance subgroups and uses cross-fitting for valid inference. \citet{abadie2024} show that exploiting first-stage heterogeneity across observable subgroups can substantially reduce the MSE of IV estimators, and that naive sample selection based on first-stage $t$-statistics causes over-rejection---a problem that our correction formula and the temporal sample split are designed to address. Principal stratification \citep{frangakis2002} handles latent compliance but typically requires parametric distributional assumptions. The contribution here is that when the experiment has temporal structure, compliance can be inferred nonparametrically and the bias from misclassification has a closed-form correction, because the false positive rate is pinned down by the test size $\alpha$.

\bibliographystyle{apalike}
\bibliography{references}

\end{document}
